% Generated draft: requires mathematical and source review.
\documentclass[11pt]{article}
\usepackage[margin=25mm]{geometry}
\usepackage{fontspec}
\setmainfont{Latin Modern Roman}
\usepackage{amsmath,amssymb,mathtools,unicode-math}
\setmathfont{Latin Modern Math}
\usepackage{enumitem,needspace,xurl,hyperref}
\setlength{\parindent}{0pt}
\setlength{\parskip}{5pt}
\setlength{\emergencystretch}{3em}
\newcommand{\N}{\mathbb{N}}
\newcommand{\Z}{\mathbb{Z}}
\newcommand{\Q}{\mathbb{Q}}
\newcommand{\R}{\mathbb{R}}
\newcommand{\C}{\mathbb{C}}
\newcommand{\F}{\mathbb{F}}
\newcommand{\A}{\mathbb{A}}
\newcommand{\PP}{\mathbb{P}}
\newcommand{\E}{\mathbb{E}}
\newcommand{\eps}{\varepsilon}
\newcommand{\dd}{\,\mathrm{d}}
\begin{document}
\section*{UnsolvedMath ID: 30000408}
\textbf{Source problem code:} catalogue entry
\subsection{Definitions and mathematical statement}
Let \(K\) be a number field and \((F,\partial_F)\) the one-variable differential field over \(K\) in the source. In the usual arithmetic-function-field setup, \(F/K\) is finitely generated of transcendence degree one and \(\partial_F\) is a nonzero derivation killing \(K\); the extracted statement does not explicitly give these model hypotheses. A differential module \((M,\partial_M)\) has finite rank \(r\) over \(F\) and an additive operator satisfying
\[
 \partial_M(fm)=\partial_F(f)m+f\partial_M(m).
\]
A full system of algebraic solutions means that there exists a finite algebraic extension \(L/F\), with its extended characteristic-zero derivation, such that \(L\otimes_F M\) has an \(L\)-basis \(e_1,\ldots,e_r\) with \(\partial(e_i)=0\). It does not mean just one nonzero algebraic horizontal vector.

To define reduction, one must spread the field/curve and connection data over an arithmetic model and exclude its finitely many bad primes. For a good nonarchimedean prime \(v\) of \(K\), let \(p_v\) be its residue characteristic and \(\nabla_v\) the reduced connection. Its \(p_v\)-curvature at a derivation \(\delta\) of the reduced function field is
\[
 \psi_v(\delta)=(\nabla_{v,\delta})^{p_v}-\nabla_{v,\delta^{[p_v]}},
\]
where \(\delta^{[p_v]}=\delta^{p_v}\) is again a derivation. Trivial \(p_v\)-curvature means \(\psi_v(\delta)=0\) for every such \(\delta\). This is vanishing, not merely nilpotence. In a one-derivation presentation the second term must not be dropped unless the reduced base derivation has zero \(p_v\)-th iterate; no special coordinate is assumed.

The Grothendieck--Katz \(p\)-curvature conjecture in the supplied formulation asks for the equivalence
\[
 M\text{ has a full system of algebraic solutions}
 \quad\Longleftrightarrow\quad
 \psi_v=0\text{ for all but finitely many good primes }v.
\]
Both directions are retained. Under the customary arithmetic convention, ``almost all'' means outside a finite set of nonarchimedean places after choosing a good model, not just infinitely many primes or a density-one subsequence. The source writes rational primes \(p\) without specifying the places above them; this convention needs confirmation in its precise setup. The same missing model data determine the admissible reductions, rather than permitting a fresh choice at each prime to force vanishing.

The literature triage checked 2026-08-21 identifies the question as the general Grothendieck--Katz conjecture, reports important proved cases, and leaves the general local--global assertion open. The exact special-case hypotheses are absent. The open label and verification note preserving conditions while normalizing delimiters remain supplied metadata, not independently verified current conclusions.
\subsection{Short English statement}
Does a differential module become trivial with a full horizontal basis after a finite algebraic extension exactly when its reductions have zero \(p\)-curvature at all but finitely many primes? The complete-basis condition, both directions, true vanishing and finite-exception quantifier are preserved. The arithmetic model needed for reduction is not supplied.
\subsection{Sources}
\begin{itemize}
\item Local UnsolvedMath registry, permanent ID 30000408, source code catalogue entry, ``Grothendieck--Katz \(p\)-Curvature Conjecture,'' including original and cleaned equivalence and verification metadata.
\item \emph{The Arithmetic of Fields (2006)}, Oberwolfach Reports, pp.\ 317--378, DOI 10.4171/owr/2006/06: \url{https://doi.org/10.4171/owr/2006/06}.
\item Supplied reference ``The p-curvature conjecture for the non-abelian Gauss--Manin connection'': \url{https://arxiv.org/abs/1912.05757}.
\item Literature triage checked 2026-08-21; local record updated 2026-09-16. No external source has been visited or independently checked here.
\end{itemize}
\end{document}
